\chapter{Einleitung}

Die Erforschung von Zell-Zell-Interaktionen in Bioreaktorsystemen erfordert
wiederholbare, quantifizierende Beobachtung über Stunden bis Tage. Manuelle
mikroskopische Untersuchung ist zeitaufwendig, subjektiv und schlecht
reproduzierbar. Insbesondere die Kalzifizierung von Osteoblasten
(hFOB\,1.19) und deren Wechselwirkung mit endothelialen Zellen (hMEC\,1)
auf HAP-X-Chips stellt hohe Anforderungen an zeitliche Auflösung und
Konstanz der Beobachtungsbedingungen.

\section{Motivation}

Die biologische Fragestellung erfordert:
\begin{itemize}
    \item Langzeitbeobachtung (24--72\,h) bei 37\,\textdegree C
    \item Auflösung im Sub-Mikrometer-Bereich (Zellmorphologie, Kalziumablagerungen)
    \item Automatisierte, reproduzierbare Bildaufnahme
    \item Echtzeit-KI-Inferenz zur Zellerkennung und -verfolgung
\end{itemize}

\section{Zielsetzung}

Ziel dieser Arbeit ist die Entwicklung eines kostengünstigen, vollständig
automatisierten Mikroskopiesystems, das:
\begin{enumerate}
    \item Open-Source-Hardware (OpenFlexure-Mikroskop) und -Software einsetzt,
    \item eine angepasste Optik für die Pi HQ Camera (IMX477) mit 125--150\,mm
          Tubusobjektiv nutzt,
    \item segmentierte KI-Inferenz (YOLOv8s-seg) auf dem Hailo-8-NPU in Echtzeit
          ausführt,
    \item Zellverfolgung (ByteTrack) und morphologische Analyse ermöglicht, und
    \item über Ansible reproduzierbar bereitgestellt wird.
\end{enumerate}

\section{Aufbau der Arbeit}

Kapitel~\ref{chap:grundlagen} gibt einen Überblick über die biologischen und
technischen Grundlagen. Kapitel~\ref{chap:methoden} beschreibt den Aufbau des
Systems. Kapitel~\ref{chap:ergebnisse} präsentiert die experimentellen
Ergebnisse. Kapitel~\ref{chap:diskussion} diskutiert die Befunde und
Kapitel~\ref{chap:zusammenfassung} fasst die Arbeit zusammen.